\documentclass[11pt]{article}
\usepackage[T1]{fontenc}
\usepackage{lmodern,microtype,amsmath,amssymb,amsthm,mathtools}
\usepackage[margin=1in]{geometry}
\usepackage{booktabs,longtable,array,tikz-cd,enumitem,xurl}
\usepackage[hidelinks]{hyperref}
\hypersetup{pdftitle={Recursion categories and their coalgebraic representations},pdfsubject={Catalogue and classification problems over sets of integers, version 3}}
\setlength{\emergencystretch}{2em}
\newtheorem{theorem}{Theorem}[section]
\newtheorem{proposition}[theorem]{Proposition}
\newtheorem{corollary}[theorem]{Corollary}
\theoremstyle{definition}
\newtheorem{definition}[theorem]{Definition}
\newtheorem{problem}[theorem]{Problem}
\newcommand{\N}{\omega}
\newcommand{\Baire}{\mathbb B}
\newcommand{\Z}{\mathbb Z}
\newcommand{\R}{\mathcal R}
\newcommand{\C}{\mathcal C}
\newcommand{\Pfn}{\mathbf{Pfn}}
\newcommand{\Set}{\mathbf{Set}}
\newcommand{\Par}{\operatorname{Par}}
\newcommand{\Coalg}{\operatorname{Coalg}}
\newcommand{\Tot}{\operatorname{Tot}}
\newcommand{\id}{\operatorname{id}}
\newcommand{\dom}{\operatorname{dom}}
\newcommand{\Rep}{\operatorname{Rep}}
\title{Recursion categories and their coalgebraic representations\\\large A catalogue and classification problems over sets of integers}
\author{}
\date{Version 3 --- 21 October 2026}
\begin{document}
\maketitle
\begin{abstract}
We catalogue classical, arithmetical, analytical, admissible, syntactic,
and operator examples of recursion categories, together with equivariant
and existence constructions. The operator categories of Montagna and
Sorbi give direct connections with effective descriptive set theory and
Wadge reducibility. We organize coalgebraic representations and their
comparisons in a 2-category, classify representations with constant
coalgebra targets up to equivalence, and construct representations using
binary observations of subsets of the integers. Countable recursion
categories admit countable coalgebraic representations. Faithfulness
requires retaining distinctions between states with equal observable
behavior. The catalogue separates the underlying recursion category,
its presentation, and the structure preserved by a representation.
Pointed Moore coalgebras also present strictly causal operators on
streams; we give their countable and computable realizations and relate
their equality to behavioral equivalence.
\end{abstract}

\section{The objects to be classified}

A catalogue of recursion categories must specify both arrows and their
equality. These determine which computations the category admits and
which computations it identifies. A coalgebraic representation adds
states, observations, and transitions; its comparison theory should
retain the partial-map structure that supports universal computation.

\subsection{Sets, sequences, and numerical codings}\label{sec:codes}
We work in ZFC and write $\omega=\{0,1,2,\ldots\}$ and
$2=\{0,1\}$. Baire space is $\Baire=\omega^\omega$ and Cantor
space is $2^\omega$, both with the product topology of their discrete
coordinate spaces. A finite prescription of coordinates defines a
basic clopen set. We identify a subset $A\subseteq\omega$ with its
characteristic function $\chi_A\in2^\omega$.

Fix the Cantor pairing bijection
\begin{equation}\label{eq:cantor-pair}
 \langle-,-\rangle:\omega^2\longrightarrow\omega,
 \qquad
 \langle m,n\rangle=\frac{(m+n)(m+n+1)}2+n.
\end{equation}
On the diagonal $m+n=s$, its values are the consecutive integers
from $s(s+1)/2$ to $(s+1)(s+2)/2-1$. These intervals partition
$\omega$, proving bijectivity and giving a computable inverse.
For each fixed arity, tuples are coded by iterating this pairing.

For a total function $f:\omega\to\omega$, define its coded graph
$\Gamma_f\subseteq\omega$ and graph characteristic function
$\chi_f:\omega\to2$ by
\begin{equation}\label{eq:graph-code}
 \Gamma_f=\{\langle m,n\rangle:f(m)=n\},\qquad
 \chi_f(\langle m,n\rangle)=
 \begin{cases}1,&f(m)=n,\\0,&f(m)\ne n.\end{cases}
\end{equation}
Thus $\chi_f=\chi_{\Gamma_f}$ is a point of Cantor space. The
graph codes of all total functions form the subspace
\[
 \mathcal G=\{x\in2^\omega:
       \text{for every }m\text{ there is exactly one }n
       \text{ with }x(\langle m,n\rangle)=1\}.
\]
This subspace is $G_\delta$, meaning a countable intersection of
open sets. Indeed, the condition of at most one $1$ in each row is
closed, and the condition of at least one is
\[
 \bigcap_{m\in\omega}\ \bigcup_{n\in\omega}
       \{x:x(\langle m,n\rangle)=1\}.
\]
Closed subsets of Cantor space are $G_\delta$, so their intersection
has the asserted form. The map $f\mapsto\chi_f$ is a homeomorphism
$\Baire\to\mathcal G$: each graph coordinate tests one value of
$f$, and the inverse image in $\mathcal G$ of the condition
$f(m)=n$ is the clopen coordinate condition
$x(\langle m,n\rangle)=1$. Both maps are computable; to recover
$f(m)$ from $\chi_f$, search row $m$ for its unique $1$.

For $f,g\in\Baire$, interleaving is the function
$(f\oplus g)(2n)=f(n)$ and $(f\oplus g)(2n+1)=g(n)$.
Its inverse projections are $P_1h(n)=h(2n)$ and
$P_2h(n)=h(2n+1)$. Write $f^+(n)=f(n+1)$ for the tail.
To code subsets of $\mathbb Z$ by subsets of $\omega$, use the
bijection $\zeta:\mathbb Z\to\omega$ given by
$\zeta(n)=2n$ for $n\geq0$ and $\zeta(n)=-2n-1$ for $n<0$.

\subsection{Partial maps and recursion categories}\label{sec:partial}
A partial function $f:X\rightharpoonup Y$ is a function
$D_f\to Y$ with $D_f\subseteq X$. Its source is $X$ and its
domain of definition is $D_f$. Its restriction $\overline f:X\rightharpoonup X$
is the identity on $D_f$. Composition $gf$ applies $f$ first;
its domain is $\{x\in D_f:f(x)\in D_g\}$.

Abstractly, a restriction category assigns $\overline f:X\to X$
to each arrow $f:X\to Y$, subject to
\[
 f\overline f=f,\qquad
 \overline f\,\overline g=\overline g\,\overline f,\qquad
 \overline{g\overline f}=\overline g\,\overline f,\qquad
 \overline g f=f\overline{gf}.
\]
The middle two equations concern arrows with the same source; the
last concerns composable $f:X\to Y$ and $g:Y\to Z$.
An arrow is \emph{total} if $\overline f=1_X$.
A \emph{restriction idempotent}, or \emph{domain}, is an arrow
$e:X\to X$ with $e=\overline e$. Domains are ordered by
$e\leq d$ when $ed=e$. For parallel arrows, the restriction order
is $f\leq g$ when $f=g\overline f$; for partial functions this
means that $g$ extends $f$.

We use the restriction-product formulation of a $P$-category
\cite[Proposition 4.5]{CL3}. A binary near product
$X\boxtimes Y$ has total projections $p_X,p_Y$ and, for every
$f:A\to X$ and $g:A\to Y$, a unique arrow
$\langle f,g\rangle:A\to X\boxtimes Y$ satisfying
\[
 p_X\langle f,g\rangle=f\overline g,\qquad
 p_Y\langle f,g\rangle=g\overline f,\qquad
 \overline{\langle f,g\rangle}=\overline f\,\overline g.
\]
Here angle brackets pair arrows; in \eqref{eq:cantor-pair} they
pair natural numbers. Set $\Delta_X=\langle1_X,1_X\rangle$ and
$f\boxtimes g=\langle fp_X,gp_Y\rangle$.
The total arrows form a category $\Tot(\C)$, in which these are
binary categorical products. No nullary product is assumed.

A $P$-category is \emph{prodominical} if it has zero arrows,
absorbing under composition, and $f\boxtimes0=0$ for every typed
pair. These equations imply $\overline0=0$. A nonzero object has
$1_X\ne0$; the category is nontrivial if it has such an object.
An \emph{isotype} is a nonempty category in which every two objects
are isomorphic. A \emph{Turing morphism} is an arrow
$\tau:W\boxtimes X\to Y$ such that, for every object $V$ and
every $f:V\boxtimes X\to Y$, there is a total $g:V\to W$ with
\[
 f=\tau(g\boxtimes1_X).
\]
A \emph{recursion category} is a prodominical $P$-category that is
an isotype and has a Turing morphism \cite[\S9]{Heller}.

Di Paola--Heller's dominical convention additionally requires
$f\boxtimes g=0$ only when $f=0$ or $g=0$ and characterizes total
arrows by $fh=0\Rightarrow h=0$ for every composable $h$
\cite[Definitions 2.0--2.3]{DPH}. The $P$-recursion categories
of Di Paola--Montagna need not be dominical. We retain this
distinction in the catalogue.

A range of $f:X\to Y$ is a domain $r:Y\to Y$ such that $rf=f$
and, for every parallel pair $u,v:Y\to Z$, the equality $uf=vf$
implies $ur=vr$. Coproducts below are categorical binary coproducts
in the partial-map category. They are distributive if the canonical map
$(A\boxtimes X)+(A\boxtimes Y)\to A\boxtimes(X+Y)$ is an
isomorphism. A functor preserves near products or coproducts when
the corresponding canonical comparison maps are isomorphisms.

Several properties distinguish the historical examples. A section of
$f:X\to Y$ is an arrow $s:Y\to X$ with $fs=\overline s$ and
$fsf=f$; \emph{choice} means that every arrow has a section.
An arrow $f:X\to Y$ is a partial monomorphism if $fu=fv$ implies
$\overline f u=\overline f v$ for every parallel $u,v:Z\to X$.
\emph{Weak choice} means that every partial monomorphism is a
partial isomorphism: it has $s:Y\to X$ with
$sf=\overline f$ and $fs=\overline s$.
An \emph{atom} is a nonzero total arrow $a:X\to Y$ such that
$ea$ is either $0$ or $a$ for every domain $e$ on $Y$.
A \emph{constant} is a total $c:X\to Y$ with $cu=cv$ for every
pair of total arrows $u,v:Z\to X$. There are \emph{enough atoms}
(respectively, constants) if precomposition by all atoms
(respectively, constants) with target $Y$ separates arrows out of
$Y$ \cite[\S\S6,8--9]{DPH}.

A nonzero domain $e$ is \emph{connected} if a disjoint decomposition
$e=e_0\vee e_1$ into domains forces one summand to be zero.
Here $\vee$ is the domain join when the indicated join exists.
The connected-domain separation condition says that $fh=gh$ for
all connected domains $h$ on the common source implies $f=g$.
This is the condition called \emph{locally connected} in the cited
recursion-category literature \cite[\S4]{DPM}.

The catalogue is organized around the following data: its objects and arrows; equality;
the near product and Turing morphism; domains and ranges; choice,
coproducts, and connectedness; and a specified coalgebraic representation.
The initial catalogue below includes the constructions of
Di Paola--Heller, Montagna, Montagna--Sorbi, and Di Paola--Montagna.

\section{The catalogue}
\subsection{Numerical and higher-recursion examples}

\paragraph{The classical category $\mathcal T$.}
Di Paola--Heller's objects are countably infinite sets with recursive
structures, drawn from a product-closed family; arrows are partial
recursive maps in those structures, with extensional equality
\cite[Theorem 4.3, pp.~601--602]{DPH}. Positive finite powers of $\N$
give a concrete equivalent presentation. Pairing gives the near product,
and a universal partial recursive function together with the parameter
theorem gives a Turing morphism. The domains are recursively enumerable
sets. Numerical search gives range sections.

\paragraph{Oracle relatives $\mathcal T_A$.}
For a fixed $A\subseteq\N$, replace partial recursive maps by partial
$A$-recursive maps on positive finite powers of $\N$. Equality remains
extensional. The oracle universal machine and parameter theorem prove
the recursion axioms by the same argument. These categories have ranges
and sections obtained by dovetailing relative to $A$.

These concrete presentations detect Turing reducibility exactly:
\begin{equation}\label{eq:oracle-order}
 \mathcal T_A\subseteq\mathcal T_B
 \quad\Longleftrightarrow\quad A\leq_T B.
\end{equation}
The inclusion means identical objects and inclusion of the actual
partial functions. One implication follows by simulating $A$-queries
using $B$. Conversely, the characteristic function of $A$ is an arrow
of $\mathcal T_A$ and hence, under the inclusion, is $B$-recursive.
Equation~\eqref{eq:oracle-order} concerns these numerical presentations;
classification under arbitrary abstract equivalence is a further question.

\paragraph{The arithmetical hierarchy.}
For $n\geq1$, Di Paola--Heller construct
$\mathcal T(\Sigma^0_n)$, whose arrows have lightface $\Sigma^0_n$
graphs in the chosen numerical structures. The domains are
$\Sigma^0_n$ sets, and the decidable domains are $\Delta^0_n$ sets.
Their enumeration argument supplies a Turing morphism
\cite[\S4.6, Example 4, p.~606]{DPH}.
For fixed numerical codings these are also the oracle categories
$\mathcal T_{0^{(n-1)}}$: a functional graph is c.e. in the indicated
jump exactly when its partial function is computable in that jump.

\paragraph{Analytical examples on integers.}
Di Paola--Heller construct $\mathcal T_\omega(\Pi^1_1)$ and state
the corresponding construction of $\mathcal T_\omega(\Sigma^1_2)$
\cite[\S4.6, Example 5, pp.~606--607]{DPH}. Arrows are partial
numerical functions whose graphs lie in the indicated lightface
pointclass. The construction of a universal functional relation uses
single-valuedness results for those pointclasses. The subscript
$\omega$ specifies that the free arguments are integers.
This is a direct entry point for effective descriptive set theory:
uniformization and universal relations supply the Turing morphism.

\paragraph{Admissible recursion.}
For admissible $\alpha$, Di Paola--Heller describe
$\mathcal T^\alpha$, using carriers that are copies of $L_\alpha$
with $\alpha$-recursive structures and partial
$\alpha$-recursive maps \cite[\S4.6, Example 7, p.~608]{DPH}.
The admissible set and its coding belong to the catalogue entry.
When $L_\alpha$ is countable, an external integer coding is available;
the coding and the internal notion of $\alpha$-recursiveness are
separate data.

\subsection{Valid syntactic categories}

\paragraph{Montagna's $S$.}
The objects are generated from $\N$ by binary products and disjoint
sums. Arrows are program indices modulo a single PA-provable universally
quantified equality of partial computations
\cite[pp.~105--108]{Montagna}. This is the case $S_{\mathrm{PA}}$
of the following construction.

\paragraph{Di Paola--Montagna's $S_T$.}
Let $T$ be a consistent recursively enumerable extension of PA. With
the same object types, indices $e,d:A\rightharpoonup B$ define the
same arrow when
\begin{equation}\label{eq:uniform}
 T\vdash\forall x\in A\,
 \bigl[(\varphi_e(x)\downarrow\leftrightarrow\varphi_d(x)\downarrow)
 \land(\varphi_e(x)\downarrow\Rightarrow
                \varphi_e(x)=\varphi_d(x))\bigr].
\end{equation}
The product and sum types use the fixed effective numerical codings.
The equality in \eqref{eq:uniform} is a congruence: substitution into
one uniform proof proves equality after precomposition, and the
arithmetized composition equations prove equality after postcomposition.
The structural identities and the parameter theorem have PA-provable
versions, so they pass to the quotient.

The resulting category is a $c^+$-$P$-recursion category, with ranges,
choice, and compatible distributive coproducts
\cite[\S2, p.~646]{DPM}. It is atomless and has insufficient constants.
If $T$ is $\Sigma^0_1$-sound, the nonzero connected domains are
exactly the provable singleton domains; otherwise there are no nonzero
connected domains. No $S_T$ is locally connected
\cite[Theorem 4.3 and Corollary 4.4, pp.~657--658]{DPM}.

The pointwise constructions $S'$ and $S'_T$ use a different equivalence
relation. Their proposed composition fails to be independent of
representatives \cite{Pointwise}; they are excluded from the catalogue.
The uniform constructions $S$ and $S_T$ remain included.

\subsection{Montagna--Sorbi's operator categories}

We use the interleaving identification $\Baire^2\cong\Baire$
from Section~\ref{sec:codes}.
An ordinary partial recursive operator can output a partial numerical
function. A partial $\Baire$-operator is defined only when that
output is total.

\paragraph{$\mathcal B$.}
This one-object category has carrier $\Baire$ and the extensional
partial recursive $\Baire$-operators as arrows
\cite[Definition 2.1, p.~1027]{MS}. Its domains are exactly the
lightface $\Pi^0_2$ subsets of Baire space. It does not have ranges for
all arrows, and weak choice fails. It has atoms but not enough atoms
\cite[Propositions 2.1--2.3 and Lemma 2.2, pp.~1027--1030]{MS}.

\paragraph{$\mathcal B_o$.}
The object is again Baire space. Arrows are all partial continuous maps
with $G_\delta$ domains, with extensional equality
\cite[Lemma 2.5 and Definition 2.2, pp.~1030--1031]{MS}.
This category contains $\mathcal B$, has enough atoms, and still lacks
ranges and weak choice \cite[Proposition 2.4]{MS}.

Both categories are recursion categories
\cite[Proposition 3.1, p.~1032]{MS}. If $(\Phi_e)$ enumerates the
partial recursive $\Baire$-operators,
a common Turing morphism is
\[
 \tau(f\oplus g)=\Phi_{f(0)}(f^+\oplus g).
\]
A total compiler prefixes an operator index to the parameter function.
The passage from $\mathcal B$ to $\mathcal B_o$ permits arbitrary
real parameters and continuous compilers.

For subsets $D,E\subseteq\Baire$, Wadge reducibility
$D\leq_W E$ means that $D=u^{-1}(E)$ for some total continuous
$u:\Baire\to\Baire$. A $G_\delta$ set is $G_\delta$-complete
if every $G_\delta$ subset of Baire space reduces to it in this sense.
Complete domains in $\mathcal B_o$ are the $G_\delta$-complete sets
under Wadge reducibility \cite[Remark 3.1, p.~1034]{MS}.
Some domains of $\mathcal B$ become complete in $\mathcal B_o$
without being complete in $\mathcal B$
\cite[Corollary 3.4, pp.~1034--1035]{MS}. Thus the comparison changes
the admitted reductions, a fact that must be recorded alongside any
coalgebraic representation.

\paragraph{The atomless category $\mathcal N$.}
Let
\[
 X=\{f\in\N^{\N}:\text{every partial recursive }\varphi_e\subseteq f
                         \text{ has finite domain}\}.
\]
For $\Phi\in\mathcal B$, let $\Phi^X$ retain precisely those inputs
in $X$ whose outputs lie in $X$. The arrows of $\mathcal N$ form the
category \emph{generated} by these restricted operators. The assignment
$\Phi\mapsto\Phi^X$ need not preserve composition. Montagna--Sorbi
prove a normal form $P_1\Phi^X$ and then establish that $\mathcal N$
is a recursion category \cite[Definitions 4.1--4.2, Lemma 4.3,
Proposition 4.1, pp.~1038--1040]{MS}. It is atomless and has diagonal
ranges. Its creativity notion relative to atoms does not imply
completeness \cite[Theorems 4.1--4.2]{MS}.

Its carrier has a useful topological description.
\begin{proposition}\label{prop:dense}
The set $X$ defining $\mathcal N$ is a dense $G_\delta$ subset of
Baire space.
\end{proposition}
\begin{proof}
For each $e$ such that $\dom\varphi_e$ is infinite, the set $C_e$ of
total extensions of $\varphi_e$ is closed. Given a basic cylinder,
choose an argument in $\dom\varphi_e$ beyond its specified coordinates
and extend the cylinder by a different value there. Hence $C_e$ is
nowhere dense. The complement of the union of these countably many
$C_e$ is $X$. The Baire category theorem proves density.
\end{proof}
This is a boldface $G_\delta$ description. Its indexing condition is
not an effective enumeration of open complements.

\subsection{Equivariant and existence constructions}

Let $G$ be a finite group and $\C$ a $+$-dominical recursion category
with equality, in the terminology of Di Paola--Heller: coproducts
have the required compatibility with the near product, and each
diagonal has a range. In $\C^G$,
write $G\otimes X$ for the free $G$-object, formed from the $G$-indexed
coproduct of $X$. The full subcategories on free objects $G\otimes X$
and on objects $(G\otimes X)+X$, with trivial action on the last
summand, are denoted $\C_1^G$ and $\C_0^G$. Both are recursion
categories \cite[Theorem 7.9.9, pp.~620--621]{DPH}.
These entries retain the finite group and the chosen base category.

Heller's existence theorem applies to uniformly generated isotypical
$B^\#$-subcategories of total maps in locally connected iteration categories
\cite[Theorem 9.2]{Heller}. His examples include set-based,
cocommutative-coalgebra, functor-category, and frame-relative
constructions \cite[\S10]{Heller}. The generating data and iteration
closure specify the resulting recursion category.

Lengyel's coalgebra construction uses partial morphisms of
$F$-coalgebras as an ambient $P$-category \cite{Coalgebras}.
The corrected relative existence theorem is as follows
\cite[Theorem 6.4]{Corrections}: let $\C,\mathcal D$ be iteration
categories, with $\mathcal D$ separated, and let
$L:\C\to\mathcal D$ be a faithful iteration functor. Every uniformly
generated isotypical $B^\#$-subcategory $\mathcal B$ of the total
maps of $\C$ has a recursion-category completion $\operatorname{Tur}(\mathcal B)$.
Here separated is the connected-domain separation condition called
locally connected in Heller's terminology.

In ZFC, if $F:\Set\to\Set$ weakly preserves pullbacks and
$\Coalg(F)$ has binary products, its partial maps with subcoalgebra
domains form a relatively separated iteration category. The relative
theorem therefore gives recursion categories from its uniformly
generated isotypical $B^\#$-subcategories of total maps
\cite[Theorem 10.1]{Corrections}. These are construction schemes:
an entry for an individual example includes the actual functor,
isotypical object, generators, and required preservation hypotheses.

\section{Representations and their 2-categorical comparisons}

For an endofunctor $F:\Set\to\Set$, an $F$-coalgebra is a map
$c:X\to FX$. A coalgebra morphism from $(X,c)$ to $(Y,d)$ is a
function $u:X\to Y$ satisfying $du=(Fu)c$; these form
$\Coalg(F)$. A subcoalgebra is a subset $D\subseteq X$ with a
coalgebra structure making its inclusion a coalgebra morphism.
A partial coalgebra morphism $(X,c)\rightharpoonup(Y,d)$ consists
of a subcoalgebra domain $D\hookrightarrow X$ and a total coalgebra
morphism $D\to Y$. Composition pulls back the second domain along
the first map. We use this construction when subcoalgebra inclusions
are stable under these pullbacks and it supplies the restriction,
zero, and binary near-product structures specified below. Write
$\mathcal A_F=\Par(\Coalg(F))$. For constant functors and the
Moore functor used in Section~4, domains and composition have this form.

A representation is a faithful functor $J:\R\to\mathcal A_F$
that preserves restrictions and zeros and has canonical total
isomorphisms
\[
 c^J_{X,Y}=\langle Jp_X,Jp_Y\rangle:
 J(X\boxtimes Y)\longrightarrow JX\boxtimes JY.
\]
Thus preservation of products means preservation through their
canonical comparison maps. Their associativity and symmetry
compatibilities follow from the product equations. One may also
require preservation of distributive binary coproducts, of ranges,
or of both; the same requirements are imposed on every functor
in a given comparison. Range preservation means
$J(\operatorname{ran}f)=\operatorname{ran}(Jf)$ whenever these
ranges are required. A Turing morphism remains universal for the
represented arrows by applying $J$ to its factorization equation.

For such functors $P,Q:\C\to\mathcal D$, a
\emph{structure-compatible natural transformation}
$\theta:P\Rightarrow Q$ has total components
$\theta_X:PX\to QX$, satisfies
$Qf\,\theta_X=\theta_Y\,Pf$ for every $f:X\to Y$, and obeys
\begin{equation}\label{eq:product-compatibility}
 c^Q_{X,Y}\theta_{X\boxtimes Y}
       =(\theta_X\boxtimes\theta_Y)c^P_{X,Y}.
\end{equation}
When coproduct preservation is required, put
$d^P_{X,Y}=[Pi_X,Pi_Y]:PX+PY\to P(X+Y)$ and require
$\theta_{X+Y}d^P_{X,Y}=d^Q_{X,Y}(\theta_X+\theta_Y)$.
These equations use the canonical comparisons, so they are stable
under vertical composition and whiskering.

To compare different catalogue entries as well as representations of
one entry, use the following construction.
\begin{definition}\label{def:global}
An object is a triple $(\R,F,J)$ as above. A 1-cell
$(\R,F,J)\to(\mathcal S,G,J')$ is a triple $(H,K,\kappa)$ consisting
of structure-preserving functors $H:\R\to\mathcal S$ and
$K:\mathcal A_F\to\mathcal A_G$, with an invertible
structure-compatible natural transformation
$\kappa:KJ\Rightarrow J'H$. A 2-cell to $(H',K',\kappa')$ is a pair
of structure-compatible natural transformations with total components,
$\alpha:H\Rightarrow H'$ and $\eta:K\Rightarrow K'$, satisfying
\[
 (J'\alpha)\kappa=\kappa'(\eta J).
\]
\end{definition}
In particular, the component
$\kappa_X:K(JX)\to J'(HX)$ is a total isomorphism, and its
naturality equation is
$J'(Hf)\kappa_X=\kappa_YK(Jf)$ for $f:X\to Y$.
The equation defining a 2-cell is an equality of transformations
$KJ\Rightarrow J'H'$. The square being compared is
\[
\begin{tikzcd}[column sep=large,row sep=large]
\R\arrow[r,"H"]\arrow[d,"J"'] & \mathcal S\arrow[d,"J'"]\\
\mathcal A_F\arrow[r,"K"'] & \mathcal A_G.
\end{tikzcd}
\]
Composition is pasting: for a second 1-cell $(H_2,K_2,\kappa_2)$,
the composite has functors $H_2H,K_2K$ and comparison
$(\kappa_2H)(K_2\kappa)$. The 2-category laws follow from those
of functors and natural transformations. Restricting to a fixed
$\R$, with $H=\id_{\R}$ and $\alpha=\id$, recovers
$\Rep(\R)$ of \cite{Comparisons}.

An equivalence in $\Rep(\R)$ includes an equivalence of the whole
ambient targets. An equivalence of the categories of represented arrows
is weaker. In fact, all faithful represented images of the same source,
with source-object tags retained, are equivalent to that source.
Classification of representations therefore requires the ambient
category or additional effective and topological data.

\subsection{Classification for constant targets}
For a set $E$, coalgebras of the constant functor $F_E(X)=E$ are
labelled sets $a:X\to E$. Their partial-map category $\Pfn_E$ has
label-preserving partial functions as arrows. Products of total
coalgebras are fibre products over $E$. A representation
$J:\R\to\Pfn_E$ has fibre functors $J_e:\R\to\Pfn$ for $e\in E$.

\begin{theorem}\label{thm:constant}
Let $J:\R\to\Pfn_E$ and $J':\R\to\Pfn_D$ be representations with
the same specified preservation requirements. They are equivalent
in $\Rep(\R)$ if and only if there is a bijection $\sigma:E\to D$
and structure-compatible natural isomorphisms
$J_e\cong J'_{\sigma(e)}$ for all $e\in E$.
\end{theorem}
\begin{proof}
A restriction equivalence $K:\Pfn_E\simeq\Pfn_D$ restricts to an
equivalence of total-map categories $\Set/E\simeq\Set/D$.
It preserves the terminal object of these total-map categories and
its subobject lattice. The latter is $\mathcal P(E)$ or
$\mathcal P(D)$. Atoms are precisely singleton labels, so $K$
induces a bijection $\sigma:E\to D$.

Every labelled set is the coproduct of its singleton labelled elements.
Preservation of these coproducts identifies $K$, naturally on total
maps, with relabelling along $\sigma$. A partial map factors as the
partial inverse of its domain inclusion followed by a total map.
Restriction preservation extends this identification to partial maps.
Thus $KJ\cong J'$ gives exactly the displayed fibre isomorphisms.
Conversely, relabelling along $\sigma$ is a structured equivalence,
and the fibre isomorphisms assemble to its required comparison with
$J'$. The inverse relabelling gives the inverse 1-cell.
\end{proof}

The label set, including labels unused by the image, is therefore an
ambient invariant. The theorem classifies the constant-target part
of $\Rep(\R)$; comparisons into larger coalgebra categories supply
additional 1-cells.

\section{Countable representations and integer-set coalgebras}
\subsection{A countable representation theorem}

\begin{proposition}\label{prop:countable}
Every nontrivial prodominical $P$-category with countably many objects
and arrows has a faithful representation in $\Pfn_E$ with $E$
countable and every represented carrier countable. It preserves
restrictions, zeros, and binary near products, and reflects totality.
If distributive binary coproducts exist, such a representation can
also preserve them.
\end{proposition}
\begin{proof}
Split restriction idempotents: the objects of the resulting category
$\mathcal K$ are $(X,e)$ with $e=\overline e$, and arrows satisfy
$f=dfe$. Its identity at $(X,e)$ is $e$, and it is countable.
Regard each original object $X$ as $(X,1_X)$.
Use its nonzero objects $A$ as labels. The $A$-fibre of $J_0X$ is
the set of total $h:A\to X$; define $J_0f(h)=fh$ precisely when
$fh$ is total.

A total composite has total inner factor, so this rule respects
composition and definedness. The identity
$\overline{fh}=\overline{\overline f h}$ proves restriction
preservation; a zero composite cannot be total on nonzero $A$.
Total maps into a restriction product are exactly pairs of total maps
with the same source, proving product preservation.
To distinguish $f$ from $g$, split $\overline f=mp$.
If $f\ne0$, $m$ is a generalized element on which $J_0f$ is
defined. Equality $J_0f=J_0g$ would give $fm=gm$, hence
$f=g\overline f$. Reversing the roles gives $g=f\overline g$,
and therefore $f=g$. The case $f=0$ follows from the reverse
argument. This proves faithfulness. Restriction preservation and
faithfulness give reflection of totality. All fibres and labels
are countable.

Suppose now that binary coproducts are distributive. The split category
$\mathcal K$ inherits these coproducts and their distributivity. A
\emph{summand domain} on $A$ is the restriction idempotent that is the
identity on one side of a decomposition $A\cong A_0+A_1$ and zero on
the other. Write $B(A)$ for the set of these idempotents.

We first record the decomposition property used below. A total map
$h:A\to X+Y$ decomposes $A$ into the inverse images of the two
summands. Indeed, put $q_0=[1_X,0]$, $q_1=[0,1_Y]$, and
$e_i=\overline{q_i h}$. Distributivity gives the total decision map
\[
 d=(p_A+p_A)\delta^{-1}\langle1_A,h\rangle:A\longrightarrow A+A,
\]
where $\delta:(A\boxtimes X)+(A\boxtimes Y)\to
A\boxtimes(X+Y)$ is the distributivity isomorphism. If
$e_i=j_i p_i$ are splittings in $\mathcal K$, then
$[j_0,j_1]:A_0+A_1\to A$ has inverse $(p_0+p_1)d$.
To check the inverse identities, let $k_i:A\to A+A$ be the
injections, $r_0=[1_A,0]$, $r_1=[0,1_A]$, and
$\nabla=[1_A,1_A]$. The product equations and naturality of
distributivity give
\[
 \nabla d=1_A,\qquad r_i d=e_i,\qquad
 (h+h)d=(i_X+i_Y)h.
\]
For the middle identity, for example,
$r_0(p_A+p_A)=p_A(1_A\boxtimes q_0)\delta$; the two sides
agree on each summand. For the last, use
$(h\boxtimes1)\langle1_A,h\rangle=\Delta h$ and check
$(p_{X+Y}+p_{X+Y})\delta^{-1}\Delta=i_X+i_Y$ on the
two injections. Applying $q_0+q_1$ gives
$(q_0h+q_1h)d=h$. Since $h$ is total, the restriction axioms
give $(e_0+e_1)d=d$. They also give $de_i=k_i e_i$ from
$\overline{r_i}=k_i r_i$ and $r_i d=e_i$.
Consequently
\[
 [j_0,j_1](p_0+p_1)d=[e_0,e_1]d
       =\nabla(e_0+e_1)d=1_A.
\]
The reverse composite is the identity on each summand because
$dj_i=k_i j_i$ and $p_i j_i=1_{A_i}$.
This proves the decomposition used in the germ construction
of \cite[Appendix A]{Representation}.
Pulling one decomposition back along the injections of another gives
a common four-part decomposition. Consequently $B(A)$ is a Boolean
algebra: meet is composition, complement exchanges the two summands,
and join takes their union. A summand of a summand is again a summand.

Use every pair $(A,U)$, where $A$ is nonzero and $U$ is an
ultrafilter on $B(A)$, as a label. An arrow $h:A\to X$ is \emph{admissible at $U$}
if $e\leq\overline h$ for some $e\in U$. Two admissible arrows
represent the same germ when
\[
 [h]_U=[k]_U
 \quad\Longleftrightarrow\quad
 \text{there is }e\in U\text{ with }
 e\leq\overline h\,\overline k\text{ and }he=ke.
\]
Intersections of supporting members of $U$ prove transitivity, so
this is an equivalence relation. Let the $(A,U)$-fibre of $JX$ be
the set of these germs, and define the partial action by
\[
 Jf([h]_U)=[fh]_U
 \qquad\text{if }fh\text{ is admissible at }U.
\]
If $h$ and $k$ agree on a member of $U$, intersect that member with
one supporting $fh$. On the intersection, $fh$ and $fk$ are defined
and equal. Thus both definedness and value are independent of the
representative. Admissibility of $gfh$ implies that of $fh$, which
proves functoriality. The identity
$\overline{fh}=\overline{\overline f h}$ proves preservation of
restrictions, and a proper ultrafilter contains no zero member, so
$J0$ is nowhere defined.

At each label, pairing gives a bijection
\[
 J(X\boxtimes Y)_{(A,U)}
 \longrightarrow JX_{(A,U)}\times JY_{(A,U)},
 \qquad [h]_U\longmapsto([p_Xh]_U,[p_Yh]_U).
\]
Its inverse pairs representatives after intersecting their supporting
members. These bijections commute with postcomposition and preserve
projections and diagonals. For a germ into $X+Y$, restrict a
representative to a supporting summand $e\in U$ and apply the
decomposition property. Its two inverse-image summands have disjoint
domains whose join is $e$. Exactly one belongs to $U$, so the germ
factors through exactly one coproduct injection. The two resulting
images are disjoint: applying $[1_X,0]$ to an equality between them
would make a supported arrow zero. Hence
$J(X+Y)\cong JX\amalg JY$, with the images of the injections as
coproduct injections.

To prove faithfulness, suppose $Jf=Jg$ and split
$\overline f=mp$, with $m:A\to X$ total. If $f\ne0$, then $A$
is nonzero and $fm$ is total. At every ultrafilter $U$ on $B(A)$,
equality of the two represented partial functions supplies some
$e\in U$ such that
\[
 e\leq\overline{gm},\qquad fme=gme.
\]
These agreement domains have a finite subfamily whose join is $1_A$.
Otherwise their complements would have the finite intersection
property and would extend to an ultrafilter meeting none of them.
Refine the finite cover to a disjoint summand partition. Its
injections are jointly epimorphic, so $fm=gm$, and hence
$f=g\overline f$. This equation also holds when $f=0$.
Interchanging $f$ and $g$ gives equality by antisymmetry of the
restriction order. Thus $J$ is faithful.

Each fibre is countable because its germs are represented by arrows
of $\mathcal K$. For each unequal parallel pair $f,g$, choose one
label at which $Jf$ and $Jg$ differ, and retain just these labels.
There are countably many pairs. The restricted functor remains
faithful, and all the preservation arguments above apply separately
at each retained label. Its carriers are countable unions of
countable fibres, as required.
\end{proof}

Consequently the classical, fixed-oracle, lightface arithmetical and
analytical, syntactic, and recursive-operator entries with countable
arrow sets have countable representations. This conclusion concerns
cardinality. An effective presentation must additionally specify
computable structure and represented arrow maps.

\subsection{Binary observations of integer subsets}
Take $M(Y)=2\times Y^{\{+,-\}}$. For $A\subseteq\Z$, let $Q_A$
have carrier $\Z$, output $\chi_A$, and transitions $n\mapsto n+1$
and $n\mapsto n-1$. Call $A$ aperiodic if $A+k=A$ implies $k=0$.

\begin{proposition}\label{prop:shift}
Every coalgebra morphism $Q_A\to Q_B$ is a translation $n\mapsto n+k$
with $B=A+k$. For aperiodic $A,B$, their translation classes therefore
classify the $Q_A$ up to isomorphism, and between two such coalgebras there is at
most one morphism. Each $Q_A$ has no proper nonempty subcoalgebra.
\end{proposition}
\begin{proof}
Commutation with both transitions forces $h(n)=h(0)+n$ for every
integer $n$. Preservation of output says
$\chi_B(n+h(0))=\chi_A(n)$, exactly the displayed translation
condition. Two allowable translations differ by a period of $A$.
A nonempty transition-closed subset of $\Z$ contains every integer.
\end{proof}

The same argument classifies all nonempty transitive binary systems
with inverse transitions. Put
$\operatorname{Stab}(A)=\{k\in\Z:A+k=A\}$.
\begin{proposition}\label{prop:transitive}
A nonempty binary Moore coalgebra with mutually inverse transitions
and no proper nonempty subcoalgebra is isomorphic to
\[
 Q(H,A)=(\Z/H,o_A,\delta_+,\delta_-),
 \qquad H\leq\operatorname{Stab}(A),
\]
where $A\subseteq\Z$ and $o_A:\Z/H\to2$ is given by $o_A(n+H)=\chi_A(n)$,
and the transitions add or subtract $1$. The containment
$H\leq\operatorname{Stab}(A)$ makes this output well defined.
Two such coalgebras are
isomorphic exactly when their subgroups $H$ agree and their observation
sets are integer translates. Moreover,
\[
 \operatorname{End}(Q(H,A))=\operatorname{Aut}(Q(H,A))
       \cong\operatorname{Stab}(A)/H.
\]
\end{proposition}
\begin{proof}
The inverse transitions define a $\Z$-action. Its orbits are
subcoalgebras, so the action is transitive. Choose a state $q$;
its stabilizer is $H$, and let $A$ consist of the integers $n$ for
which the output at $n\cdot q$ is $1$. This identifies the carrier
with $\Z/H$ and implies $H\leq\operatorname{Stab}(A)$.
Changing the chosen state translates $A$ and leaves $H$ unchanged.
An equivariant map between two transitive $\Z$-sets is determined
by the image of the zero coset. For an isomorphism the stabilizers
must agree; output preservation gives exactly the translation
condition. For an endomorphism the allowable cosets are
$\operatorname{Stab}(A)/H$, and every such translation is invertible.
\end{proof}

Thus the rigid systems in this class have $H=\operatorname{Stab}(A)$.
They are the shift orbits of integer subsets, with finite cycles for
periodic subsets and infinite orbits for aperiodic subsets.

For a countable label set choose $E\subseteq\N$ and put
$A_e=\{0,e+1\}$. These sets are aperiodic and pairwise nontranslate.
For a labelled set $(X,a)$ define $\Phi(X)=X\times\Z$, with
\[
 o(x,n)=\chi_{A_{a(x)}}(n),\qquad
 \delta_+(x,n)=(x,n+1),\quad \delta_-(x,n)=(x,n-1).
\]
For a label-preserving partial function $f:D\to Y$, put
$\Phi(f)=f\times\id_{\Z}$ with domain $D\times\Z$.

\begin{theorem}\label{thm:integer}
The functor
$\Phi:\Pfn_E\to\Par(\Coalg(M))$ is fully faithful and preserves
restrictions, zeros, coproducts, and all coalgebra products indexed
by nonempty sets on total maps, hence binary near products. Composing it with the
representation in Proposition~\ref{prop:countable} gives a faithful
countable coalgebraic representation with binary observations.
\end{theorem}
\begin{proof}
A subcoalgebra of $X\times\Z$ is $D\times\Z$ because both
translations are available. A morphism on each copy of $\Z$ lands
in one target copy; Proposition~\ref{prop:shift} and the chosen
markers force preservation of its label and integer phase. It is
therefore exactly $f\times\id$ for a labelled partial function $f$.
This proves full faithfulness and the assertions about domains and zeros.

For a family $(X_i,a_i)_{i\in I}$ with $I\ne\varnothing$, put
\[
 \prod_E X_i=\{(x_i)\in\prod_{i\in I}X_i:
                     a_i(x_i)\text{ is independent of }i\}.
\]
Its label is this common value. The coalgebra product of the
$\Phi(X_i)$ is $(\prod_E X_i)\times\Z$. To verify the universal property, take
any cone of total coalgebra maps into the factors. At a state of its
vertex, all cone components have the same output after every input
word. In the image objects these outputs determine the marker label
and integer phase, so both agree in every coordinate. The component
maps thus pair uniquely into the stated object and commute with the
transitions. Coproducts are disjoint unions with inherited structure.
The product description also respects partial domains.
\end{proof}

The marker systems are uniformly computable in $e$. The generalized
element or germ construction may nevertheless give noneffective
represented arrows. For an effective refinement, computability of
these partial actions is a separate mathematical requirement.

\begin{corollary}\label{cor:reals}
A nontrivial prodominical $P$-category with at most $2^{\aleph_0}$
objects and arrows has a faithful representation by coalgebras for
$M(Y)=2\times Y^{\{+,-\}}$ whose states can be coded by subsets of
$\N$. Restrictions, zeros, and binary near products are preserved;
distributive binary coproducts can also be preserved when present.
\end{corollary}
\begin{proof}
The proof of Proposition~\ref{prop:countable}, with an infinite
cardinal $\kappa$ in place of countability, bounds labels and carriers
by $\kappa$. In the germ construction there are at most $\kappa$
parallel arrow pairs to separate, and every fibre has size at most
$\kappa$. Apply this with $\kappa=2^{\aleph_0}$.
For the marker construction choose distinct subsets $A_e\subseteq\N$
containing $0$. Their minimum is $0$, so each is aperiodic as a subset
of $\Z$, and distinct such sets cannot be translates. There are
$2^{\aleph_0}$ choices. The proof of Theorem~\ref{thm:integer}
applies to this family. Each resulting carrier has cardinal at most
$2^{\aleph_0}$ and therefore injects into $\mathcal P(\N)$;
transport its coalgebra structure along the injection.
\end{proof}

This applies to $\mathcal B_o$ as well as the countable-arrow examples.
There are at most continuum many $G_\delta$ domains, and a continuous
map from each such domain into Baire space is determined by its values
on a countable dense subset. The resulting set-theoretic representation
leaves the topological preservation problem in Section~6 to be solved.

\section{Descriptive set theory and computation}
\subsection{Sets as behaviors and as codes}

There are several useful ways for sets of integers to enter a
representation. A state may be an oracle, a function coded by its graph,
or a complete coalgebra presentation coded as a real. Alternatively,
an integer set may be the observable behavior of a state. These
interpretations impose different conditions on morphisms.

The graph coding of Section~\ref{sec:codes} identifies the carriers
of $\mathcal B$, $\mathcal B_o$, and $\mathcal N$ with subspaces
of $\mathcal G\subseteq2^\omega$. It also makes their operators
partial maps between sets of graph codes, by conjugating with
$f\mapsto\chi_f$. To obtain a coalgebraic representation one must
choose observations and transitions for which the admitted arrows
are coalgebra morphisms.

\subsection{Binary and positive observations}
The coalgebra
\[
 \mathcal P(\N)\longrightarrow2\times\mathcal P(\N),\qquad
 S\longmapsto\bigl(\chi_S(0),\{n:n+1\in S\}\bigr)
\]
is final for $2\times-$ \cite[Corollary 2.3.6]{Jacobs}.
For $c=(o,\delta):X\to2\times X$, its unique behavior map is
$b(x)=\{n:o(\delta^n x)=1\}$. Every $S$ has a countable
realization on $\N$, with successor transition and output $\chi_S$.
A realization with computable initial state and total computable
binary output and transition exists exactly when $S$ is recursive:
iteration computes each membership bit, and the successor construction
proves the converse. A finite-state realization exists exactly when
the characteristic sequence is ultimately periodic.

Positive information uses Sierpi\'nski space $\mathbb S$, whose open
sets are $\varnothing,\{1\},\{0,1\}$. The Scott topology on
$\mathcal P(\N)$ has basic opens $\{S:D\subseteq S\}$ for finite
$D\subseteq\N$, and identifies this space with $\mathbb S^{\N}$.
The same head and tail make it final for $\mathbb S\times-$ in
topological spaces: the displayed behavior map is continuous because
each positive membership test has open inverse image.
More explicitly, an enumeration name for $S\subseteq\omega$ is a
sequence $p\in\omega^\omega$ satisfying
\[
 S=\{n:(\exists k)\ p(k)=n+1\}.
\]
The value $0$ supplies a blank entry, and repetitions are allowed.
The set has a computable enumeration name exactly when it is c.e.
By contrast, the characteristic-function representation names $S$
by its binary sequence $\chi_S$; that name is computable exactly
when $S$ is recursive \cite{Scott,Pauly}.

Enumeration operators supply another direct connection. A relation
$W\subseteq\operatorname{Fin}(\N)\times\N$ defines
\[
 \Psi_W(S)=\{n:(D,n)\in W\text{ for some finite }D\subseteq S\}.
\]
These are Scott-continuous maps; conversely, a Scott-continuous map
is determined by its values on finite sets and has such a form.
For the converse, if $n\in\Psi(S)$, Scott openness of the inverse
image of $\{T:n\in T\}$ supplies a finite $D\subseteq S$ with
$n\in\Psi(D)$. Thus one can take
$W=\{(D,n):n\in\Psi(D)\}$.
To make effectiveness precise, code a finite set $D$ by
$\ulcorner D\urcorner=\sum_{k\in D}2^k$, and code its rule for
$n$ by $\langle\ulcorner D\urcorner,n\rangle$.
If this set of rule codes is c.e., $\Psi_W$ is an enumeration
operator: while enumerating $S$, output $n$ whenever all members
of some enumerated rule's finite premise $D$ have appeared.
The information supplied by observations thus distinguishes decision
from semidecision.

\subsection{Why observable behavior alone is insufficient}
\begin{proposition}\label{prop:multiplicity}
Let $\R$ be a nontrivial recursion category and
$J:\R\to\Par(\Coalg(F))$ a faithful representation preserving
totality. Suppose $F$ has a final coalgebra $B$. For every object $X$
of $\R$, the behavior morphism $JX\to B$ is not monic in
$\Coalg(F)$.
\end{proposition}
\begin{proof}
An object $Q$ with monic map to a final object is subterminal: for
any $u,v:P\to Q$, their composites into the final object agree,
so $u=v$. Choose a Turing morphism $\tau:W\boxtimes U\to V$.
The isotype supplies a total isomorphism
$h:W\boxtimes U\to V$. It is nonzero since every object of a
nontrivial isotype is nonzero. Apply the Turing property to $h$ and zero.
It supplies distinct total endomorphisms $g_h,g_0:W\to W$, since
equal parameter maps would give $h=0$. Their faithful images are
distinct total coalgebra endomorphisms of $JW$. Hence $JW$ is not
subterminal. Every $JX$ is isomorphic to $JW$, which proves the claim.
\end{proof}

The copies in Theorem~\ref{thm:integer} retain exactly this necessary
kind of multiplicity: two states can have the same label and phase
but different copy indices. Forgetting those indices can identify
distinct represented arrows. Domains are also relevant: a partial
map descends along a behavioral quotient only when its domain is
saturated under behavioral equivalence \cite[Proposition 3.3]{Comparisons}.

\subsection{Equivalence relations on coalgebra codes}\label{sec:moore-codes}
Fix a finite or countable input alphabet $A$, with a chosen coding
of its finite words. The space of pointed binary Moore coalgebras
on carrier $\omega$ is
\[
 \mathcal C_A=2^\omega\times(\omega^\omega)^A\times\omega.
\]
A code $(o,(\delta_a)_{a\in A},x)$ specifies the output function
$o:\omega\to2$, one total transition
$\delta_a:\omega\to\omega$ for each $a$, and an initial state
$x\in\omega$. All coordinates have their product topologies, with
$\omega$ discrete. The behavior of this pointed system is
\[
 L_x=\{w\in A^*:o(\delta_wx)=1\}\in\mathcal P(A^*).
\]
For $w=a_1\cdots a_k$, set
$\delta_w=\delta_{a_k}\cdots\delta_{a_1}$, so the letters are
read from left to right; $\delta_\varepsilon=1$ for the empty word.
The map $\beta:\mathcal C_A\to2^{A^*}$ sending a code to
$\chi_{L_x}$ is continuous. To determine one output coordinate
$w$, one reads only the initial state and the finitely many table
entries encountered while following $w$.

A bisimulation between two such systems is a relation on their
carriers whose related states have equal outputs and whose
$a$-successors are related for every $a\in A$. Related states
have equal languages, by induction on word length. Conversely,
the relation of having equal languages is a bisimulation, since
equal languages have equal empty-word outputs and equal residual
languages after each letter. Hence pointed bisimilarity is
\[
 E_A=(\beta\times\beta)^{-1}(\Delta_{2^{A^*}}).
\]
It is closed because $\beta$ is continuous and the diagonal of
$2^{A^*}$ is closed. An equivalence relation on a standard Borel
space is \emph{smooth} if a Borel map into $2^\omega$ reduces it
to equality. Coding $A^*$ by integers makes $\beta$ such a
reduction, so $E_A$ is smooth. For finite $A$, the behaviors
realized by finite-state systems are exactly the regular languages
\cite{RuttenAutomata}.

For countable nondeterministic labelled transition systems,
S\'anchez Terraf proves that bisimilarity is
$\boldsymbol\Sigma^1_1$-complete \cite{Terraf}. Thus the chosen
transition functor changes the descriptive complexity of behavioral
classification. Isomorphism of presentations and equivalence of
representations are separate equivalence relations.

\subsection{Moore coalgebras and higher-type operators}\label{sec:higher-type}

The finite-word behavior of a pointed Moore coalgebra also determines
an operator on infinite input streams. Let $A$ be a nonempty countable
alphabet. For a stream $\alpha\in A^\omega$, write
$\alpha\upharpoonright n=\alpha(0)\cdots\alpha(n-1)$ for its
prefix of length $n$, including the empty prefix when $n=0$.
A pointed Moore coalgebra with output $o:X\to2$, transitions
$\delta_a:X\to X$, and initial state $x\in X$ induces
\begin{equation}\label{eq:moore-stream}
 T_x:A^\omega\longrightarrow2^\omega,
 \qquad T_x(\alpha)(n)=o(\delta_{\alpha\upharpoonright n}x).
\end{equation}
The zeroth output is $o(x)$, emitted before the first input is read.
We call an operator $T:A^\omega\to2^\omega$ \emph{strictly causal}
if, for all $\alpha,\gamma\in A^\omega$ and $n\in\omega$,
\begin{equation}\label{eq:strict-causality}
 \alpha\upharpoonright n=\gamma\upharpoonright n
 \quad\Longrightarrow\quad T(\alpha)(n)=T(\gamma)(n).
\end{equation}
Thus the output at time $n$ depends only on the preceding $n$ inputs.
Every $T_x$ is strictly causal, and every strictly causal operator is
continuous for the product topologies: any finite set of output
coordinates is determined by a sufficiently long finite input prefix.

When $A=\omega$, the displayed operator has type
$(\omega\to\omega)\to(\omega\to2)$. With numerical and Boolean
values at type level zero and streams at level one, this is an operator
of type level two. Each coordinate $\alpha\mapsto T_x(\alpha)(n)$
is a Boolean-valued functional of type two. The coalgebra's output and
transition maps give a state-based presentation of this operator.

For the effective statements, take $A=\omega$ or
$A=\{0,\ldots,k-1\}$ with $k\geq1$. An operator
$T:A^\omega\to2^\omega$ is \emph{Type-2 computable} if one oracle
algorithm, given $n$ and oracle access to $\alpha$, halts with output
$T(\alpha)(n)$ for every $\alpha$ and $n$. Each such computation
reads finitely many coordinates of its oracle. A \emph{computable
Moore presentation} has carrier $\omega$, a computable output
$o:\omega\to2$, a jointly computable transition map
$(a,s)\mapsto\delta_a(s)$ from $A\times\omega$ to $\omega$,
and a specified initial state in $\omega$.

\begin{proposition}\label{prop:moore-realization}
For every nonempty countable alphabet $A$, the strictly causal
operators $A^\omega\to2^\omega$ are exactly the operators induced
by pointed Moore coalgebras with countable state sets. If $A=\omega$
or $A=\{0,\ldots,k-1\}$ with $k\geq1$, a strictly causal operator
is Type-2 computable if and only if it has a computable Moore
presentation.
\end{proposition}
\begin{proof}
Strict causality of a Moore operator follows from
\eqref{eq:moore-stream}. Conversely, let $T$ be strictly causal
and choose $a_0\in A$. Use the countable state set $A^*$, initial
state $\varepsilon$, and transitions and outputs
\[
 \delta_a(s)=sa,\qquad o(s)=T(sa_0^\omega)(|s|).
\]
Here $sa_0^\omega$ is the infinite stream beginning with $s$ and
continuing with $a_0$ forever. Strict causality makes this output
independent of the chosen tail. After reading
$\alpha\upharpoonright n$, the state is exactly that prefix, and
its output is $T(\alpha)(n)$.

A computable Moore presentation computes $T_x(\alpha)(n)$ by
reading the first $n$ inputs, applying the corresponding transitions,
and evaluating the output. For the converse, give $A^*$ a computable
bijection with $\omega$ and use the preceding prefix construction.
Such a bijection is obtained by length-lexicographic enumeration for
finite $A$; for $A=\omega$, enumerate words by increasing
$|s|+\sum_{i<|s|}s(i)$, with lexicographic order within each finite
group. Concatenation is computable in either coding. If $T$ is
Type-2 computable, the value $T(sa_0^\omega)(|s|)$ is computable
uniformly in $s$, since its oracle stream is uniformly computable
from $s$. Thus the output and joint transition maps are computable.
\end{proof}

The interpretation preserves exactly the behavioral equality of
\S\ref{sec:moore-codes}. For two pointed systems with languages
$L_x$ and $L_y$,
\begin{equation}\label{eq:moore-operator-equality}
 T_x=T_y\quad\Longleftrightarrow\quad L_x=L_y.
\end{equation}
Equality of languages gives equality of every output in
\eqref{eq:moore-stream}. Conversely, extend any word $w$ to
$wa_0^\omega$ and inspect output $|w|$ to recover its membership
in the language. Hence the equivalence relation $E_A$ also classifies
these pointed systems by equality of their induced stream operators.

The code-to-behavior map $\beta:\mathcal C_A\to2^{A^*}$ has a
uniform effective interpretation as well. For the effective alphabets
above, name a code by its initial state and oracle access to the
tables $o(s)$ and $\delta_a(s)$; the pairing function
\eqref{eq:cantor-pair} combines the table indices into numerical
coordinates. Given a word $w$, an oracle algorithm follows its
finitely many transitions and queries the final output. It therefore
computes $\beta(c)(w)$ uniformly from the name of $c$, even when
the tables themselves are noncomputable. With these representations,
$\beta$ is a Type-2 computable map between spaces of codes.

The timing convention distinguishes Moore operators from the full
class of continuous stream operators. For example, the left shift
$S:2^\omega\to2^\omega$, defined by $S(\alpha)(n)=\alpha(n+1)$,
is computable and continuous but fails \eqref{eq:strict-causality}.
An operator is \emph{causal} if its output at time $n$ depends only
on the input prefix of length $n+1$. Deleting the zeroth output of
a Moore operator gives a causal operator. Mealy coalgebras for
$X\mapsto(B\times X)^A$ produce an output with each input;
their final coalgebra consists of all causal functions
$A^\omega\to B^\omega$ \cite[Section 2.2, Theorem 2.3]{Mealy}.

General continuous stream operators can read a variable finite
number of inputs before emitting each output. Garner studies such
stream processors and gives an extensional coalgebraic
characterization of the continuous functions
$A^\omega\to B^\omega$ \cite[Theorem 1.1]{Garner}. For
$A=B=\omega$, these are operators on Baire space. The total
Type-2 computable operators and total continuous operators are,
respectively, the total arrows of Montagna--Sorbi's $\mathcal B$
and $\mathcal B_o$. Extending these presentations to their partial
arrows requires accounting for the $\Pi^0_2$ and $G_\delta$
domains described in \S2.3. This is part of the first classification
problem below.

\section{Specific classification problems}

The catalogue suggests three concrete comparisons.
\begin{enumerate}[leftmargin=*,itemsep=6pt]
\item For $\mathcal B\subseteq\mathcal B_o$, construct compatible
coalgebraic representations that preserve the distinction between
effective and continuous reductions. Determine which representations
recover the $\Pi^0_2$ and $G_\delta$ domain structures and the
completeness change identified by Montagna--Sorbi.
\item For $S_T$, compare the countable generalized-element and germ
representations as $T$ varies. Determine which coalgebraic invariants
detect provable equality, ranges, and the connected-domain distinction
of Di Paola--Montagna's Theorem 4.3.
\item For analytical and admissible entries, identify the uniformization
and coding needed for effective coalgebraic representations. Compare
the resulting 1-cells with numerical oracle inclusions and with
pointclass-preserving transformations.
\end{enumerate}

For the shift coalgebras of Proposition~\ref{prop:shift}, the first
isomorphism classification is explicit: integer translation of the
observation set. The next problem is to classify representations
assembled from these systems, retaining multiplicities, admissible
domains, and the represented arrow family. Theorem~\ref{thm:constant}
provides the constant-target comparison with which to begin.

\begin{thebibliography}{99}
\small
\setlength{\itemsep}{3pt}
\bibitem{CL3} J.~R.~B. Cockett and S.~Lack,
\emph{Restriction categories III: colimits, partial limits and extensivity},
Mathematical Structures in Computer Science \textbf{17} (2007), 775--817.
\url{https://arxiv.org/abs/math/0610500}.
\bibitem{DPH} R.~A. Di Paola and A.~Heller,
\emph{Dominical categories: recursion theory without elements},
Journal of Symbolic Logic \textbf{52} (1987), 594--635.
\href{https://doi.org/10.1017/S0022481200029649}{doi:10.1017/S0022481200029649}.
\bibitem{Montagna} F.~Montagna,
\emph{``Pathologies'' in two syntactic categories of partial maps},
Notre Dame Journal of Formal Logic \textbf{30} (1989), 105--116.
\href{https://doi.org/10.1305/ndjfl/1093634998}{doi:10.1305/ndjfl/1093634998}.
\bibitem{MS} F.~Montagna and A.~Sorbi,
\emph{Creativeness and completeness in recursion categories of partial recursive operators},
Journal of Symbolic Logic \textbf{54} (1989), 1023--1041.
\href{https://doi.org/10.2307/2274763}{doi:10.2307/2274763}.
\bibitem{DPM} R.~A. Di Paola and F.~Montagna,
\emph{Some properties of the syntactic $p$-recursion categories generated
by consistent, recursively enumerable extensions of Peano arithmetic},
Journal of Symbolic Logic \textbf{56} (1991), 643--660.
\href{https://doi.org/10.2307/2274707}{doi:10.2307/2274707}.
\bibitem{Heller} A.~Heller, \emph{An existence theorem for recursion categories},
Journal of Symbolic Logic \textbf{55} (1990), 1252--1268.
\href{https://doi.org/10.2307/2274486}{doi:10.2307/2274486}.
\bibitem{Coalgebras} F.~Lengyel, \emph{Recursion categories of coalgebras},
2001. \url{https://arxiv.org/abs/math/0105256}.
\bibitem{Lengyel} F.~Lengyel, \emph{More existence theorems for recursion categories},
Annals of Pure and Applied Logic \textbf{125} (2004), 1--41.
\href{https://doi.org/10.1016/S0168-0072(03)00060-5}{doi:10.1016/S0168-0072(03)00060-5}.
\bibitem{Corrections} F.~Lengyel,
\emph{Corrections to More existence theorems for recursion categories},
version 10, 2026. \url{https://github.com/flengyel/errata/tree/main/lengyel2004}.
\bibitem{Pointwise} F.~Lengyel,
\emph{Pointwise provable equality and the failure of composition},
2026. \url{https://arxiv.org/abs/2609.25556}.
\bibitem{Representation}
\emph{Coalgebraic representations of prodominical recursion categories},
manuscript, version 2, 19 October 2026.
\url{https://github.com/flengyel/Categorical_Incompleteness/blob/main/Coalgebraic_Representation_v2.tex}.
\bibitem{Comparisons} \emph{A 2-category of coalgebraic representations},
manuscript, 20 October 2026.
\url{https://github.com/flengyel/Categorical_Incompleteness/blob/main/Coalgebra_Representation_Comparisons_v1.tex}.
\bibitem{Jacobs} B.~Jacobs, \emph{Introduction to coalgebra: towards mathematics of states and observation},
online manuscript, Chapter 2.
\url{https://www.cs.ru.nl/B.Jacobs/CLG/JacobsCoalgebraIntro.pdf}.
\bibitem{RuttenAutomata} J.~J.~M.~M. Rutten,
\emph{Automata and coinduction (an exercise in coalgebra)},
CONCUR 1998, Lecture Notes in Computer Science \textbf{1466}, 194--218.
\url{https://ir.cwi.nl/pub/2114/2114D.pdf}.
\bibitem{Mealy} H.~H. Hansen, D.~de Oliveira Costa, and J.~J.~M.~M. Rutten,
\emph{Synthesis of Mealy machines using derivatives},
CWI report SEN-R0605, May 2006.
\url{https://ir.cwi.nl/pub/10858/10858D.pdf}.
\bibitem{Garner} R.~Garner, \emph{Stream processors and comodels},
Logical Methods in Computer Science \textbf{19}(1) (2023), article 2, 1--24.
\href{https://doi.org/10.46298/LMCS-19(1:2)2023}{doi:10.46298/LMCS-19(1:2)2023}.
\url{https://arxiv.org/abs/2106.05473}.
\bibitem{Scott} D.~Scott, \emph{Data types as lattices},
SIAM Journal on Computing \textbf{5} (1976), 522--587.
\url{https://www.cs.ox.ac.uk/files/3287/PRG05.pdf}.
\bibitem{Pauly} A.~Pauly,
\emph{On the topological aspects of the theory of represented spaces},
\url{https://arxiv.org/abs/1204.3763}.
\bibitem{Terraf} P.~S\'anchez Terraf, \emph{Bisimilarity is not Borel},
Mathematical Structures in Computer Science \textbf{27} (2017), 1265--1284.
\url{https://arxiv.org/abs/1211.0967}.
\end{thebibliography}
\end{document}
